\documentclass[10pt]{article}
\usepackage[T1]{fontenc}
\usepackage[utf8]{inputenc}
\usepackage[french]{babel}
\usepackage{amsmath,amssymb,mathrsfs,mathpazo}
\usepackage{geometry}
\geometry{a4paper,margin=2cm}
\usepackage{graphicx,array,tabularx,tabularray,colortbl,multirow,multicol,booktabs,enumitem,pifont,ulem,fancybox,diagbox,cancel}
\usepackage{pstricks,pstricks-add,pst-plot,pst-tree,pst-node,pst-text,pst-func,pst-math,pst-bspline,pst-3dplot,pst-eucl,pst-solides3d}
\let\pscalculate\undefined
\usepackage{pst-bezier}
\usepackage{tikz,pgfplots,tkz-tab}
\usetikzlibrary{arrows,arrows.meta,patterns,calc,positioning,trees,decorations.markings,shapes.geometric,shapes.misc,shapes.arrows,intersections,decorations.pathmorphing,decorations.pathreplacing,backgrounds,shadows,babel}
\pgfplotsset{compat=1.18}
\usepgfplotslibrary{fillbetween}
\usepackage[np]{numprint}
\usepackage{listings,makecell,pas-tableur,twoopt}
\let\barre\undefined
\usepackage[pstricks]{tablvar}
\definecolor{darkgreen}{RGB}{0,100,0}
\providecommand{\red}{\color{red}}
\providecommand{\blue}{\color{blue}}
\providecommand{\green}{\color{green}}
\providecommand{\black}{\color{black}}
\providecommand{\white}{\color{white}}
\providecommand{\gray}{\color{gray}}
\providecommand{\cyan}{\color{cyan}}
\providecommand{\magenta}{\color{magenta}}
\providecommand{\yellow}{\color{yellow}}
\providecommand{\violet}{\color{violet}}
\providecommand{\orange}{\color{orange}}
\providecommand{\euro}{\texteuro}
\providecommand{\EUR}{\texteuro}
\providecommand{\ieme}{\textsuperscript{e}}
\providecommand{\degres}{\ensuremath{^\circ}}
\providecommand{\degree}{\ensuremath{^\circ}}
\providecommand{\vv}[1]{\overrightarrow{#1}}
\providecommand{\Coord}[1]{#1}
\providecommand{\decosix}{}
\providecommand{\starredbullet}{\ensuremath{\bullet}}
\pagestyle{empty}
\setlength{\parindent}{0pt}
\usepackage{adjustbox}\DeclareUnicodeCharacter{2236}{:}\DeclareUnicodeCharacter{2190}{\ensuremath{\leftarrow}}
\providecommand{\R}{\mathbb{R}}
\providecommand{\N}{\mathbb{N}}
\providecommand{\Z}{\mathbb{Z}}
\providecommand{\Q}{\mathbb{Q}}
\providecommand{\D}{\mathbb{D}}
\providecommand{\C}{\mathbb{C}}
\providecommand{\e}{{\mathrm{e}}}
\providecommand{\vect}[1]{\overrightarrow{#1}}
\providecommand{\vectt}[1]{\overrightarrow{\mathrm{#1}}}
\providecommand{\vvt}[1]{\overrightarrow{\mathrm{#1}}}
\providecommand{\np}[2][]{{#2~#1}}
\providecommand{\vv}[1]{{\overrightarrow{#1}}}
\renewcommand{\d}{{\mathrm{d}}}
\providecommand{\degres}{{^\circ}}
\providecommand{\degre}{{^\circ}}
\providecommand{\blue}{\color{blue}}
\providecommand{\red}{\color{red}}
\providecommand{\enskip}{\quad}
\providecommand{\llbracket}{\lbrack\!\lbrack}
\providecommand{\rrbracket}{\rbrack\!\rbrack}
\providecommand{\euro}{{\text{€}}}
\providecommand{\no}{{\text{n°}}}
\providecommand{\interval}[2]{{[#1;#2]}}
\providecommand{\intervFO}[2]{{[#1;#2[}}
\providecommand{\dotfill}{\ldots}

\providecommand{\ds}{\displaystyle}
\providecommand{\textding}[1]{#1}
\providecommand{\Oij}{\(\mathrm{O};\vec{\imath},\vec{\jmath}\)}
\providecommand{\Oijk}{\(\mathrm{O};\vec{\imath},\vec{\jmath},\vec{k}\)}
\providecommand{\Ouv}{\(\mathrm{O};\vec{u},\vec{v}\)}
\providecommand{\intervFF}[2]{\ensuremath{[#1;#2]}}
\renewcommand{\euro}{\ifmmode\text{\texteuro}\else\texteuro\fi}
\DeclareUnicodeCharacter{00B0}{\ifmmode^{\circ}\else\textdegree\fi}
\DeclareUnicodeCharacter{20AC}{\ifmmode\text{\texteuro}\else\texteuro\fi}
% Énoncés : Éducation nationale. Transcription et corrigés : APMEP, Denis Vergès et contributeurs. Découpage et classement : CoopMaths, Sébastien Lozano, Éric Elter et contributeurs. Adaptation de la présentation : Maths Explorer.
% CC BY-SA : https://coopmaths.fr/www/about/
% https://coopmaths.fr/alea/static/bac/2026/tex/bac_2026_06_sujet2_metropole_1.tex
\begin{document}
\noindent\textit{APMEP / CoopMaths. CC BY-SA. Présentation adaptée par Maths Explorer.}\par\medskip
\begin{minipage}{0.45\linewidth}
		Dans le repère orthonormé \Oijk{}, on considère les points :

		A$(2 ; 1 ; 1)$, B$(3 ;-2 ; 0)$, C$(0 ;-1 ; 1)$ et D$(0 ; 0 ; 2)$.
	\end{minipage}
	\hfill
	\begin{tikzpicture}[x=(15:15mm),y=(135:10mm),z=(90:16mm),>=stealth,baseline={(current bounding box.center)}]
		\draw[darkgray] (-0.25,0,0)--(1.47,0,0) (2.9,0,0)--(4,0,0)
			(0,-3,0)--(0,2,0)
			(0,0,-1.5)--(0,0,3);
		\draw[darkgray,dashed] (1.47,0,0)--(2.9,0,0) ;
		\draw[->,line width=1.2pt] (0,0,0)--(1,0,0) node[pos=0.6,below]{$\vect{\imath}$} ;
		\draw[->,line width=1.2pt] (0,0,0)--(0,1,0) node[pos=0.9,below]{$\vect{\jmath}$} ;
		\draw[->,line width=1.2pt] (0,0,0)--(0,0,1) node[pos=0.7,left]{$\vect{k}$} ;
		\node at(0,0,0) [below left] {O};
		\foreach \coord/\nom/\pos in {{(2,1,1)}/A/{above},{(3,-2,0)}/B/{above right},{(0,-1,1)}/C/{below left},{(0,0,2)}/D/{above left}}
		{\fill[shift={\coord}] (0,0,0) circle (0.1pt) node[\pos] {\nom} ;}
		\draw[fill opacity=0.3,fill=lightgray!50] (2,1,1)--(3,-2,0)--(0,-1,1)--(0,0,2)--cycle;
		\draw (3,-2,0)--(0,0,2);
		\draw[dashed] (2,1,1)--(0,-1,1) ;
	\end{tikzpicture}


	\begin{enumerate}
		\item Montrer que les points A, B et C définissent un plan.
		\item \begin{enumerate}
			\item Montrer que le vecteur $\vect{n} \begin{pmatrix} 1 \\ -1 \\ 4 \end{pmatrix} $ est normal au plan (ABC).

			\item En déduire qu'une équation cartésienne du plan (ABC) est :
			\[x-y+4 z-5=0.\]
		\end{enumerate}

		\item Vérifier qu'une représentation paramétrique de la droite $\Delta$ passant par D et orthogonale au plan (ABC) est donnée par :
		\[\begin{cases}
			x=t \\
			y=-t \\
			z=2+4 t
		\end{cases} \quad \text { avec } t \in \R.\]
		\item Montrer que le projeté orthogonal de D sur le plan (ABC) est le point H$\left(-\dfrac{1}{6} ; \dfrac{1}{6} ; \dfrac{4}{3}\right)$.
		\item \begin{enumerate}
			\item Montrer que le triangle ABC est isocèle en B.

			\item Montrer que l'aire du triangle ABC est $3 \sqrt{2}$.
		\end{enumerate}

		\item \begin{enumerate}
			\item En déduire que le volume du tétraèdre ABCD est égal à 1.

			\emph{On rappelle que le volume d'un tétraèdre est donné par:} \quad $V=\dfrac{1}{3} \mathcal{B} h$,\quad \emph{ où } $\mathcal{B}$ \emph{est l'aire d'une base et} $h$ \emph{la hauteur relative à cette base.}

			\item On admet que la distance du point A au plan (BCD) est égale à $\sqrt{2}$.

			En déduire l'aire du triangle BCD.
		\end{enumerate}

		\item Dans cette question, on note $\mathrm{D}_{k}$ le point ayant pour coordonnées $(0 ; 0 ; k)$, où $k$ est un réel.
		\begin{enumerate}
			\item Déterminer la valeur de $k$ pour laquelle les points A, B, C et $\mathrm{D}_{k}$ sont coplanaires.

			Quel est alors le projeté orthogonal du point D$_{k}$ sur le plan (ABC)?

			\item Peut-on trouver une valeur de $k$ telle que A soit le projeté orthogonal de $\mathrm{D}_{k}$ sur le plan (ABC) ?
		\end{enumerate}
	\end{enumerate}

\end{document}
